\begin{circuitikz}[scale=0.9, transform shape]
  % Terminal de entrada vi
  \draw (0,0) node[ground] {}
    to[open, v^>=$v_i$] (0,2.5)
    to[short, o-*] (1.5,2.5) coordinate (in_node);

  % Resistencia RG = R1 // R2 a masa de CA
  \draw (in_node) to[R, l=$R_G$, -] (1.5,0) node[ground] {};

  % Compuerta aislada (circuito abierto)
  \draw (in_node) to[short] (3.0,2.5) coordinate (gate_pt);
  \draw (3.0,0) node[ground] {} coordinate (source_pt);
  \draw (gate_pt) to[open, v^>=$v_{gs}$] (source_pt);

  % Fuente controlada gm * vgs orientada hacia abajo (desde drenador hacia masa)
  \draw (5.5,2.5) coordinate (drain_pt)
    to[cI, l_=$g_m v_{gs}$] (5.5,0) node[ground] {};

  % Resistencia RD y RL en paralelo a masa
  \draw (drain_pt) to[short, *-*] (7.0,2.5) coordinate (rd_node)
    to[R, l=$R_D$, -] (7.0,0) node[ground] {};

  \draw (rd_node) to[short, -*] (8.5,2.5) coordinate (rl_node)
    to[R, l=$R_L$, -] (8.5,0) node[ground] {};

  % Terminal de salida sobre el drenador
  \draw (rl_node) to[short, -o] (9.8,2.5) node[above] {$v_L$} coordinate (out_node);
  \draw (9.8,0) node[ground] {} to[open, v_>=$v_L$] (out_node);
\end{circuitikz}
